\documentclass[11pt,a4paper]{article}
\usepackage[T1]{fontenc}
\usepackage{lmodern,amsmath,amssymb,amsthm,mathtools,booktabs,microtype}
\usepackage[margin=26mm]{geometry}
\usepackage[colorlinks=true,linkcolor=blue,urlcolor=blue]{hyperref}
\setlength{\emergencystretch}{2em}
\newtheorem{theorem}{Theorem}[section]
\newtheorem{lemma}[theorem]{Lemma}
\newtheorem{proposition}[theorem]{Proposition}
\newtheorem{corollary}[theorem]{Corollary}
\newcommand{\Q}{\mathbb Q}\newcommand{\Z}{\mathbb Z}
\title{Integral Lucas labels from a finite quadratic choice}
\author{}\date{15 November 2026}
\begin{document}\maketitle
\begin{abstract}
An integral congruence on one coordinate of a norm-one equation forces
exponential growth, while monicity makes the other coordinate integral
without a separate integrality condition. We determine the sharp lower
bound, give positive witnesses with prescribed finite congruence data,
and record an equivalent single square equation of degree six and
coefficient height sixteen. A rational parametrization also shows why
arbitrarily accurate local approximations to the integral labels do not
give the same conclusion. All results in this paper are unconditional.
\end{abstract}

\section{The norm-one equation}
For $a\in\Z$ define the Lucas sequence
\[
 U_0(a)=0,\qquad U_1(a)=1,\qquad
 U_{j+1}(a)=aU_j(a)-U_{j-1}(a)\quad(j\ge1).
\]
For $A\ge3$ these numbers are positive and strictly increasing for
$j\ge1$. The recurrence gives
\begin{equation}\label{eq:identities}
 U_{j+1}(a)^2-aU_{j+1}(a)U_j(a)+U_j(a)^2=1,
 \qquad U_j(-a)=(-1)^{j-1}U_j(a)
\end{equation}
for $j\ge1$. The first identity also holds for $j=0$.

\begin{lemma}\label{lem:classification}
Let $a,t\in\Z$, $|a|\ge3$, and let $w\in\Q$ satisfy
\begin{equation}\label{eq:norm}
                         t^2-atw+w^2=1.
\end{equation}
Then $w\in\Z$. If $t\ne0$, there are $j\ge1$ and
$\varepsilon\in\{1,-1\}$ such that
\[
 t=\varepsilon U_j(a),\qquad
 w\in\{\varepsilon U_{j-1}(a),\varepsilon U_{j+1}(a)\}.
\]
Conversely all these displayed pairs satisfy~\eqref{eq:norm}.
For $A\ge3$ and $j\ge1$,
\begin{equation}\label{eq:elementary-growth}
                              U_j(A)\ge(A-1)^{j-1}.
\end{equation}
\end{lemma}
\begin{proof}
The number $w$ is a rational root of the monic integer polynomial
$W^2-atW+t^2-1$, and hence is an integer. For an elementary
classification first suppose $a=A\ge3$. If $tw\le0$, the equation
implies $t^2+w^2\le1$, giving only the zero-coordinate solutions.
Otherwise change the signs of both coordinates and interchange them
as necessary to obtain $t\ge w>0$. Equality $t=w$ would give
$(2-A)t^2=1$, which is impossible. The other root in the first
coordinate is
\[
                    t'=Aw-t=\frac{w^2-1}{t}.
\]
It is an integer with $0\le t'<w$. Thus replacing $(t,w)$ by
$(w,t')$ preserves the equation and strictly decreases the larger
positive coordinate. The process terminates at $(1,0)$. Reversing
these steps gives the consecutive pairs $(U_{j+1}(A),U_j(A))$.
Their interchanges and simultaneous negatives give all solutions.
For negative $a$, replace $w$ by $-w$ and use the second identity
in~\eqref{eq:identities}. This proves the displayed classification.
The converse is the first identity in~\eqref{eq:identities}.
Finally $U_2(A)=A\ge A-1$, and the recurrence and monotonicity give
$U_{j+1}(A)\ge(A-1)U_j(A)$ for $j\ge1$.
\end{proof}

\section{A congruence forcing a large index}

\begin{theorem}\label{thm:compression}
Let $a,n\in\Z$ satisfy
\[
 |a|\ge3,\qquad c=|a-2|\ge3,\qquad n\ne0,
 \qquad t=(a-2)n+1.
\]
Suppose a rational $w$ satisfies~\eqref{eq:norm}. Then $w\in\Z$ and
\begin{equation}\label{eq:sharp}
 |t|\ge U_{c-1}(|a|)\ge(|a|-1)^{c-2},\qquad
 |n|\ge\frac{U_{c-1}(|a|)-1}{c}.
\end{equation}
In particular,
\begin{equation}\label{eq:log-growth}
 \log|n|\ge(c-2)\log(|a|-1)-\log(2c).
\end{equation}
The first bound in~\eqref{eq:sharp} is sharp for every allowed $a$.
If $a\ge5$ and $n>0$, the sharper sharp bound is
\begin{equation}\label{eq:positive}
 t\ge U_{c+1}(a),\qquad
 n\ge\frac{U_{c+1}(a)-1}{c},\qquad c=a-2.
\end{equation}
\end{theorem}
\begin{proof}
Since $t\equiv1\pmod c$, one has $t\ne0$. The possibility $t=1$
would give $n=0$, while $t=-1$ would give $c\mid2$. Thus
$t\notin\{-1,0,1\}$. Lemma~\ref{lem:classification} gives
$t=\varepsilon U_j(a)$ with $j\ge2$.
As $a\equiv2\pmod c$, induction in the recurrence gives
\[
                              U_j(a)\equiv j\pmod c.
\]
Consequently $\varepsilon j\equiv1\pmod c$. Among integers $j\ge2$
in either of the two allowed residue classes, the least is $c-1$.
Monotonicity of $U_j(|a|)$ proves the first inequality in~\eqref{eq:sharp}.
The second follows from $|t|\le c|n|+1$. Since
$U_{c-1}(|a|)\ge2$, replacing $U_{c-1}(|a|)-1$ by
$U_{c-1}(|a|)/2$ proves~\eqref{eq:log-growth}.

Sharpness is supplied by
\[
 t=-U_{c-1}(a),\qquad w=-U_{c-2}(a),\qquad
 n=\frac{-U_{c-1}(a)-1}{a-2}.
\]
The congruence just proved makes $n$ integral, and $|t|>1$ makes it
nonzero. Identity~\eqref{eq:identities} proves the norm equation.
For $a\ge5,n>0$ one has $t>1$, so $\varepsilon=1$ and
$j\equiv1\pmod c$. Its least possible value is $c+1$.
Taking $t=U_{c+1}(a)$ and $w=U_c(a)$ proves both the sharper bound
and its sharpness.
\end{proof}

The exclusions on $a$ cannot simply be dropped. For $a=4,n=-1$,
one has $t=-1,w=0$; for $a=3,n=-2$ the same pair occurs.
Both satisfy the norm equation without any growth.

\begin{corollary}[A finite square condition]\label{cor:square}
For integers $a,n$ satisfying the hypotheses of
Theorem~\ref{thm:compression}, the following are equivalent:
\begin{enumerate}
\item There is $w\in\Q$ satisfying~\eqref{eq:norm}, with
      $t=(a-2)n+1$.
\item There is $z\in\Q$ such that
\begin{equation}\label{eq:square}
          z^2=(a^2-4)((a-2)n+1)^2+4.
\end{equation}
\end{enumerate}
Every such $z$ is an integer. Equation~\eqref{eq:square} has total
degree six and coefficient height sixteen after expansion and primitive
integer normalization. Its auxiliary coordinate is a finite choice
over the two retained integral coordinates $a,n$.
\end{corollary}
\begin{proof}
Completing the square gives
\[
 (2w-at)^2=(a^2-4)t^2+4.
\]
Conversely $w=(at+z)/2$ solves the norm equation. Its integrality
follows from Lemma~\ref{lem:classification}; alternatively, a rational
square root of an integer is integral, and reducing the discriminant
modulo four shows that $at$ and $z$ have the same parity.
The right side of~\eqref{eq:square} expands as
\[
 (a^4-4a^3+16a-16)n^2
 +(2a^3-4a^2-8a+16)n+a^2.
\]
The coefficient of $z^2$ is one, and the claimed degree and height
follow directly.
\end{proof}

\section{Positive witnesses and finite congruence conditions}

\begin{proposition}\label{prop:progressions}
Let $a\ge5$ be an integer, put $c=a-2$, and define
\[
 n_0=\frac{U_{c+1}(a)-1}{c},\qquad w_0=U_c(a).
\]
For every positive integer $M$ there are arbitrarily large positive
integers $n\equiv n_0\pmod M$ and integers $w\equiv w_0\pmod M$
satisfying~\eqref{eq:norm} with $t=cn+1$.
If $a\equiv1\pmod6$, then $n_0$ and $w_0$ are odd.
\end{proposition}
\begin{proof}
The matrix
\[
                    B_a=\begin{pmatrix}a&-1\\1&0\end{pmatrix}
\]
is invertible modulo $cM$, since its determinant is one. Let $T>0$
be a multiple of its order in the finite group
$\operatorname{GL}_2(\Z/cM\Z)$ and also a multiple of $c$.
For $j_k=c+1+kT$, the recurrence gives
\[
 U_{j_k}(a)\equiv U_{c+1}(a)\pmod{cM},\qquad
 U_{j_k-1}(a)\equiv U_c(a)\pmod{cM}.
\]
Set $n_k=(U_{j_k}(a)-1)/c$ and $w_k=U_{j_k-1}(a)$.
These are positive integers, have the asserted residues, and satisfy
the norm equation by~\eqref{eq:identities}. Their sizes tend to infinity
by~\eqref{eq:elementary-growth}.
If $a\equiv1\pmod6$, reduction of the recurrence modulo two has the
period $0,1,1$. Thus $U_{a-1}(a)$ is even and $U_{a-2}(a)$ is odd.
Since $c$ is odd, $n_0=(U_{a-1}(a)-1)/c$ is odd as well.
\end{proof}

\section{Local approximation and exact integral labels}
The integrality assumptions on $a,n$ in
Theorem~\ref{thm:compression} are essential. Replacing them by rational
numbers arbitrarily close to integers at finitely many primes is a
different condition, even when the norm equation holds exactly.

\begin{proposition}\label{prop:local}
For every integer $r\ge1$ there are rational numbers $n_r,w_r$ such
that, with $t_r=3n_r+1$,
\[
 t_r^2-5t_rw_r+w_r^2=1,\qquad
 v_{19}(n_r-1)\ge r,\qquad 0<n_r<1.
\]
In contrast the exact integral labels $a=5,n=1$ admit no rational
solution of the norm equation.
\end{proposition}
\begin{proof}
The integer $85$ has the simple square root $3$ modulo $19$.
Lift it to an integer $q_r$ satisfying
$q_r^2\equiv85\pmod{19^r}$ and $q_r\equiv3\pmod{19}$.
Choose $s_r\ge19$ with
\[
                       3s_r\equiv10+q_r\pmod{19^r}.
\]
Put
\begin{equation}\label{eq:param}
 D_r=s_r^2-5s_r+1,\qquad
 t_r=\frac{s_r^2-1}{D_r},\qquad
 w_r=\frac{s_r(5s_r-2)}{D_r},\qquad
 n_r=\frac{t_r-1}{3}.
\end{equation}
Substitution proves the norm equation identically in $s_r$.
Modulo $19$, one has $s_r\equiv17$ and $D_r\equiv15$,
so every denominator in~\eqref{eq:param} is a $19$-adic unit.
Moreover
\[
 3D_r(t_r-4)=-9s_r^2+60s_r-15
                  =85-(3s_r-10)^2.
\]
It follows that $v_{19}(t_r-4)\ge r$, and hence
$v_{19}(n_r-1)\ge r$. Since $s_r\ge19$, one has $D_r>0$,
$t_r>1$ from $5s_r-2>0$, and $t_r<4$ from
$3s_r^2-20s_r+5>0$. Thus $0<n_r<1$.

For the integral target $a=5,n=1$, the value of $t$ is four and
the discriminant is $340=4\cdot85$, which is not a rational square.
The equivalence in Corollary~\ref{cor:square} proves the last assertion.
\end{proof}

The square condition in Corollary~\ref{cor:square} therefore removes
a separate integral coordinate only when the other two labels are
already exact integers. Proposition~\ref{prop:local} rules out deriving
that replacement from local approximation alone. An application to
elliptic index comparisons must separately identify the retained
arguments with their integral indices.

\section*{Exact verification}
The standard-library verifier
\nolinkurl{computation/integer_labels/pell_compression_v40/check.py}
checks both polynomial identities by independent coefficient arithmetic,
the discriminant expansion, positive and sharp witnesses for both signs
of $a$, the parity assertion, the excluded small cases, finite recurrence
periods, and explicit local approximants. It also exhausts a stated
finite rectangle of integral labels using integer square roots.
These computations check the specifications and edge cases; the proofs
above establish the infinite statements. No arithmetic hypothesis from
a defining-polynomial construction is used.
\end{document}
